\begin{problem}{알파벳 블록}{standard input}{standard output}{1 second}{1024 megabytes}

S, C, O, N 4종류의 알파벳 블록으로 문자열을 만들며 놀던 도현이는 두 가지 사실을 알아냈다.

\begin{itemize}
    \item 알파벳 O 블록 1개를 반으로 자르면 알파벳 C 블록 2개로 사용할 수 있다. 알파벳 C 블록 2개를 합쳐서 알파벳 O 블록 1개로 사용하는 것도 가능하다.
    \item 알파벳 S 블록 1개를 뒤집으면 알파벳 N 블록 1개로 사용할 수 있다. 알파벳 N 블록 1개를 뒤집어서 알파벳 S 블록 1개로 사용하는 것도 가능하다.
\end{itemize}

놀이를 마친 도현이는 블록들을 상자에 넣어 가져가려 한다. 상자에는 ``SCON'' 문자열을 넣을 수 있는 공간과 ``SCCC'' 문자열을 넣을 수 있는 공간이 하나씩 존재하며, 도현이는 두 문자열이 모두 들어간 상자만을 가져갈 수 있다.

각 알파벳 블록의 개수가 주어졌을 때, 도현이가 가져갈 수 있는 상자의 최댓값을 구해보자.


\InputFile
첫째 줄에 알파벳 블록의 개수 $b_S$, $b_C$, $b_O$, $b_N$이 공백으로 구분되어 주어진다.
$(0 ≤ b_S, b_C, b_O, b_N ≤ 1,000,000)$

\OutputFile
첫째 줄에 도현이가 가져갈 수 있는 상자 개수의 최댓값을 출력한다.

\Examples

\begin{example}
\exmpfile{example.01}{example.01.a}%
\exmpfile{example.02}{example.02.a}%
\exmpfile{example.03}{example.03.a}%
\end{example}

\end{problem}

